% Trigonometrie – bank/trigonometrie/ (zusammenbau v0.4, Heft)
\einheitenkopf[t5]{Trigonometrie}
\setcounter{aufgabe}{33}

% Anlauf (ohne Prüfkennung)
\begin{aufgabe}{Seite berechnen – Anlauf}
\begin{teile}
\teil Rechter Winkel bei $B$. Schreibe an jede Seite H, G oder A – vom Winkel $\alpha$ aus.

\dreieck{(0,0)}{(5,0)}{(5,3)}{r}{t}{s}{\alpha}{}{}
\teil Rechter Winkel bei $A$. Trage den Bruch ein.

\dreieck{(0,0)}{(4,0)}{(0,3)}{x}{y}{z}{}{}{\gamma}
$\mathrm{tan}\,\gamma =$ \leerfeld
\teil Taschenrechner im Grad-Modus: $\mathrm{sin}\,23^\circ$ auf drei Dezimalen? \leerfeld
\end{teile}
\end{aufgabe}

\begin{aufgabe}{Seite berechnen – Prüfungsaufgaben}
\begin{teile}
\teil Rechter Winkel bei $C$. Trage den Bruch ein. \hfill (P10 2017 OS) \\

\dreieck{(0,0)}{(5,0)}{(1.8,2.4)}{d}{e}{f}{}{\beta}{}
$\mathrm{sin}\,\beta =$ \leerfeld
\teil Rechter Winkel bei $B$. Trage den Bruch ein. \hfill (P10 2019 OS) \\

\dreieck{(0,0)}{(5,0)}{(5,3)}{k}{m}{n}{\alpha}{}{}
$\mathrm{cos}\,\alpha =$ \leerfeld
\teil Rechter Winkel bei $A$. Welche Gleichung gilt? \hfill (P10 2018 OS) \\ \kreuz{$\mathrm{sin}\,\beta = \frac{p}{q}$} \\ \kreuz{$\mathrm{cos}\,\beta = \frac{q}{p}$} \\ \kreuz{$\mathrm{sin}\,\beta = \frac{q}{p}$}

\dreieck{(0,0)}{(4,0)}{(0,3)}{p}{q}{r}{}{\beta}{}
\teil Rechter Winkel bei $C$. Welche Gleichung gilt? \hfill (P10 2018 OS) \\ \kreuz{$\mathrm{cos}\,\alpha = \frac{w}{v}$} \\ \kreuz{$\mathrm{cos}\,\alpha = \frac{v}{w}$} \\ \kreuz{$\mathrm{sin}\,\alpha = \frac{v}{w}$}

\dreieck{(0,0)}{(5,0)}{(3.2,2.4)}{u}{v}{w}{\alpha}{}{}
\teil Rechter Winkel bei $B$. Welche Gleichung gilt? \hfill (P10 2018 OS) \\ \kreuz{$\mathrm{tan}\,\gamma = \frac{m}{k}$} \\ \kreuz{$\mathrm{tan}\,\gamma = \frac{k}{m}$} \\ \kreuz{$\mathrm{sin}\,\gamma = \frac{m}{n}$}

\dreieck{(0,0)}{(5,0)}{(5,3)}{m}{n}{k}{}{}{\gamma}
\end{teile}
\end{aufgabe}

\begin{aufgabe}{Seite berechnen – Prüfungsaufgaben – weiter}
\begin{teile}
\teil Ein Grundstück $ABCD$ hat bei $A$ einen rechten Winkel. Die Diagonale $BD$ ist $840$ m lang, der Winkel $ADB$ beträgt $47^\circ$. Zeige rechnerisch, dass $AB \approx 614$ m ist. \hfill (P10 2021 OS)
\teil Ein Segelrevier $ABCD$ hat bei $A$ einen rechten Winkel. Die Diagonale $BD$ ist $2\,300$ m lang, der Winkel $ADB$ beträgt $41^\circ$. Zeige rechnerisch, dass $AB \approx 1\,509$ m ist. \hfill (P10 2021 OS)
\teil Ein Feld $ABCD$ hat bei $A$ einen rechten Winkel. Die Diagonale $BD$ ist $780$ m lang, der Winkel $ADB$ beträgt $53^\circ$. Berechne die Länge von $AD$. \hfill (P10 2021 OS) \\ AD ≈ \leerfeld[m]
\teil Ein Waldstück $ABCD$ hat bei $A$ einen rechten Winkel. Die Diagonale $BD$ ist $1\,650$ m lang, der Winkel $ADB$ beträgt $62^\circ$. Berechne die Länge von $AD$. \hfill (P10 2021 OS) \\ AD ≈ \leerfeld[m]
\teil Im Dachgeschoss ist eine senkrechte Wand $1{,}35$ m hoch. Die Dachschräge trifft den Boden im Punkt $P$ und bildet mit ihm einen Winkel von $38{,}5^\circ$. Wie weit ist $P$ von der Wand entfernt? \hfill (P10 2017 OS) \\ \leerfeld[m]
\teil In einem Zelt ist die senkrechte Seitenwand $0{,}85$ m hoch. Die Zeltbahn läuft bis zum Boden und bildet mit ihm einen Winkel von $27{,}4^\circ$. Wie weit liegt der Hering von der Seitenwand entfernt? \hfill (P10 2017 OS) \\ \leerfeld[m]
\teil Im Dreieck $ABC$ ist die Höhe von $C$ auf $AB$ eingezeichnet, sie ist $6{,}3$ m lang. Der Winkel bei $B$ beträgt $57^\circ$. Berechne die Seite $a = BC$. \hfill (P10 2022 OS) \\

\dreieck{(0,0)}{(6,0)}{(3.5,3)}{a}{b}{c}{}{57^\circ}{}
a ≈ \leerfeld[m]
\teil Im Dreieck $ABC$ ist die Höhe von $C$ auf $AB$ eingezeichnet, sie ist $4{,}8$ cm lang. Der Winkel bei $B$ beträgt $63^\circ$. Berechne die Seite $a = BC$. \hfill (P10 2022 OS) \\

\dreieck{(0,0)}{(5,0)}{(3.8,2.4)}{a}{b}{c}{}{63^\circ}{}
a ≈ \leerfeld[cm]
\teil $\mathrm{cos}\,60^\circ = \frac{4}{x}$ – berechne $x$. \hfill (P10 2020 OS) \\ x = \leerfeld
\teil $\mathrm{tan}\,45^\circ = \frac{6}{x}$ – berechne $x$. \hfill (P10 2020 OS) \\ x = \leerfeld
\teil Rechter Winkel bei $B$. Trage den Bruch ein. \hfill (P10 2025 OS) \\

\dreieck{(0,0)}{(5,0)}{(5,3)}{e}{f}{d}{}{}{\gamma}
$\mathrm{sin}\,\gamma =$ \leerfeld
\teil Rechter Winkel bei $A$. Trage den Bruch ein. \hfill (P10 2025 OS) \\

\dreieck{(1.8,2.4)}{(0,0)}{(5,0)}{z}{x}{y}{}{\beta}{}
$\mathrm{sin}\,\beta =$ \leerfeld
\teil Rechter Winkel bei $C$. Gib $\mathrm{tan}\,\alpha$ als Bruch an. \hfill (P10 2020 OS) \\

\dreieck{(4,0)}{(0,3)}{(0,0)}{q}{p}{s}{\alpha}{}{}
$\mathrm{tan}\,\alpha =$ \leerfeld
\teil Rechter Winkel bei $A$. Gib $\mathrm{tan}\,\gamma$ als Bruch an. \hfill (P10 2020 OS) \\

\dreieck{(0,0)}{(0,3)}{(5,0)}{w}{v}{u}{}{}{\gamma}
$\mathrm{tan}\,\gamma =$ \leerfeld
\end{teile}
\end{aufgabe}

% Anlauf (ohne Prüfkennung)
\begin{aufgabe}{Winkel berechnen – Anlauf}
\begin{teile}
\teil Gegeben: Hypotenuse $7$ cm und Gegenkathete $4$ cm. Gesucht: der Winkel $\alpha$. Was ist gesucht? \\ \kreuz{Seite: Gleichung umstellen} \\ \kreuz{Winkel: Umkehrtaste}
\teil Gegenkathete $3$ cm, Hypotenuse $8$ cm – wie groß ist $\alpha$? $\alpha \approx$ \leerfeld $^\circ$
\teil Ankathete $5$ cm, Hypotenuse $7$ cm – wie groß ist $\alpha$? $\alpha \approx$ \leerfeld $^\circ$
\end{teile}
\end{aufgabe}

\begin{aufgabe}{Winkel berechnen – Prüfungsaufgaben}
\begin{teile}
\teil Eine Seilbahn führt von der Talstation $A$ zur Bergstation $B$. Das Seil ist $610$ m lang, die waagerechte Strecke unter dem Seil $460$ m. Berechne den Winkel zwischen Seil und Waagerechter. \hfill (P10 2024 OS) \\ \leerfeld $^\circ$
\teil Eine Materialseilbahn ist $245$ m lang, waagerecht überbrückt sie $190$ m. Berechne den Winkel zwischen Seil und Waagerechter. \hfill (P10 2024 OS) \\ \leerfeld $^\circ$
\teil Beim Handball steht die Spielerin $A$ $8$ m von der Torecke $C$ entfernt; der Punkt $B$ liegt so, dass bei $B$ ein rechter Winkel ist und $AB = 3$ m. Zeige rechnerisch, dass der Winkel $\alpha$ bei $A$ rund $68^\circ$ beträgt. \hfill (P10 2019 OS)
\teil Auf einem Bolzplatz ist $AB = 7$ m und $AC = 15$ m, bei $B$ liegt ein rechter Winkel. Zeige rechnerisch, dass der Winkel $\alpha$ bei $A$ rund $62^\circ$ beträgt. \hfill (P10 2019 OS)
\teil Im Parallelogramm $ABCD$ ist $BC = 26$ cm. $F$ liegt auf der Verlängerung von $AB$ über $B$ hinaus, $CF$ steht senkrecht auf $AF$ und $BF = 9$ cm. Zeige, dass der Winkel $\varepsilon = \angle CBF$ rund $70^\circ$ groß ist. \hfill (P10 2026 FOR)
\teil Im Parallelogramm $ABCD$ ist $BC = 18$ cm. $F$ liegt auf der Verlängerung von $AB$, $CF$ steht senkrecht auf $AF$ und $BF = 7$ cm. Berechne den Winkel $\varepsilon = \angle CBF$. \hfill (P10 2026 FOR) \\ $\varepsilon \approx$ \leerfeld $^\circ$
\teil Im Dreieck $ABC$ ohne rechten Winkel ist die Höhe von $C$ auf $AB$ $5{,}8$ m lang, $b = AC = 9{,}6$ m. Berechne den Winkel $\alpha$. \hfill (P10 2022 OS) \\ $\alpha \approx$ \leerfeld $^\circ$
\teil Im Dreieck $ABC$ ohne rechten Winkel ist die Höhe von $C$ auf $AB$ $3{,}9$ cm lang, $b = AC = 8{,}5$ cm. Berechne den Winkel $\alpha$. \hfill (P10 2022 OS) \\ $\alpha \approx$ \leerfeld $^\circ$
\teil Eine Rampe überbrückt waagerecht $240$ cm und steigt dabei $18$ cm. Berechne die Länge $x$ der Rampe und ihren Anstiegswinkel $\alpha$. \hfill (P10 2025 OS) \\ $x \approx$ \leerfeld[cm] , $\alpha \approx$ \leerfeld $^\circ$
\teil Eine Laderampe überbrückt waagerecht $300$ cm und steigt dabei $25$ cm. Berechne die Länge $x$ der Rampe und ihren Anstiegswinkel $\alpha$. \hfill (P10 2025 OS) \\ $x \approx$ \leerfeld[cm] , $\alpha \approx$ \leerfeld $^\circ$
\teil Die Punkte $A$, $B$, $C$ bilden ein Dreieck mit rechtem Winkel bei $A$. Berechne die Länge von $BC$ und die Größe des Winkels $\beta$ bei $B$ (Längen in Kästchen). \hfill (P10 2019 OS) \\

\begin{ksys}[xmin=-2,xmax=4,ymin=-4,ymax=3,ablesen] \punkt{-1}{-3}{A} \punkt{3}{-3}{B} \punkt{-1}{2}{C} \end{ksys}
$BC \approx$ \leerfeld , $\beta \approx$ \leerfeld $^\circ$
\teil Die Punkte $A$, $B$, $C$ bilden ein Dreieck mit rechtem Winkel bei $A$. Berechne die Länge von $BC$ und die Größe des Winkels $\beta$ bei $B$ (Längen in Kästchen). \hfill (P10 2019 OS) \\

\begin{ksys}[xmin=-1,xmax=4,ymin=-6,ymax=3,ablesen] \punkt{0}{2}{A} \punkt{3}{2}{B} \punkt{0}{-5}{C} \end{ksys}
$BC \approx$ \leerfeld , $\beta \approx$ \leerfeld $^\circ$
\end{teile}
\end{aufgabe}

% Anlauf (ohne Prüfkennung)
\begin{aufgabe}{Teildreiecke und Vermessung – Anlauf}
\begin{teile}
\teil Fahre im Trapez das rechtwinklige Teildreieck mit der Höhe $h$ links nach. Beschrifte seine Seiten mit H, G, A – vom Basiswinkel unten links aus.

\trapez[hoehe=h]{5}{3}{2.5}
\teil Gleichschenkliges Dreieck: Schenkel $7$ cm, Basiswinkel $63^\circ$. Wie lang ist die Höhe $h$ auf die Grundseite? h ≈ \leerfeld[cm]
\teil Trapez: Höhe $4$ m, Basiswinkel $62^\circ$. Wie lang ist der Schenkel $s$? s ≈ \leerfeld[m]
\end{teile}
\end{aufgabe}

\begin{aufgabe}{Teildreiecke und Vermessung – Prüfungsaufgaben}
\begin{teile}
\teil Ein Messgerät steht $1{,}4$ m über dem Boden. Die Sichtlinie zur Hügelkuppe ist $96{,}5$ m lang und steigt unter $27^\circ$ an. Wie hoch ist der Hügel? Auf der Kuppe steht ein $9$ m hoher Aussichtsplattform – wie hoch über dem Boden liegt seine Oberkante? \hfill (P10 2018 OS) \\ \leerfeld[m] ; \leerfeld[m]
\teil Ein Messgerät steht $1{,}7$ m über dem Boden. Die Sichtlinie zur Hügelkuppe ist $148$ m lang und steigt unter $24^\circ$ an. Wie hoch ist der Hügel? Auf der Kuppe steht ein $7{,}5$ m hoher Funkmast – wie hoch über dem Boden liegt seine Oberkante? \hfill (P10 2018 OS) \\ \leerfeld[m] ; \leerfeld[m]
\teil Ein rechtwinkliges Trapez hat die senkrechten Seiten $18$ cm und $7$ cm und die waagerechte Grundseite $16$ cm. $\alpha$ ist der Winkel an der oberen Ecke der $18$-cm-Seite zwischen dieser Seite und der schrägen Seite. Berechne $\alpha$. \hfill (P10 2020 OS) \\ $\alpha \approx$ \leerfeld $^\circ$
\teil Ein rechtwinkliges Trapez hat die senkrechten Seiten $13$ cm und $5$ cm und die waagerechte Grundseite $10$ cm. $\alpha$ ist der Winkel an der oberen Ecke der $13$-cm-Seite zwischen dieser Seite und der schrägen Seite. Berechne $\alpha$. \hfill (P10 2020 OS) \\ $\alpha \approx$ \leerfeld $^\circ$
\teil Im Dreieck $ABC$ liegt $F$ auf $AC$, $BF$ steht senkrecht auf $AC$. $AB = 96$ cm, der Winkel bei $A$ ist $41^\circ$, $\angle FBA = 49^\circ$, $\angle CBF = 57^\circ$. Zeige, dass $BF \approx 63{,}0$ cm und $AF \approx 72{,}5$ cm sind, und berechne $BC$. \hfill (P10 2023 OS) \\ BC ≈ \leerfeld[cm]
\teil Im Dreieck $ABC$ liegt $F$ auf $AC$, $BF$ steht senkrecht auf $AC$. $AB = 86$ cm, der Winkel bei $A$ ist $33^\circ$, $\angle FBA = 57^\circ$, $\angle CBF = 62^\circ$. Zeige, dass $BF \approx 46{,}8$ cm und $AF \approx 72{,}1$ cm sind, und berechne $BC$. \hfill (P10 2023 OS) \\ BC ≈ \leerfeld[cm]
\teil Im Parallelogramm $ABCD$ ist die Höhe $h = CF \approx 21{,}4$ cm bekannt; $F$ liegt auf der Verlängerung von $AB$. Der Winkel $CAB$ beträgt $19^\circ$. Berechne die Länge der Diagonale $AC$. \hfill (P10 2026 FOR) \\ AC ≈ \leerfeld[cm]
\teil Im Parallelogramm $ABCD$ ist die Höhe $h = CF \approx 15{,}6$ cm bekannt; $F$ liegt auf der Verlängerung von $AB$. Der Winkel $CAB$ beträgt $23^\circ$. Berechne die Länge der Diagonale $AC$. \hfill (P10 2026 FOR) \\ AC ≈ \leerfeld[cm]
\steil Im Dreieck $ABC$ ist $AD$ die Höhe auf $BC$ ($D$ zwischen $B$ und $C$), $AD = 610$ m. $\angle ACD = 59^\circ$, $\angle CAB = 72^\circ$. Berechne $BD$. \hfill (P10 2016 OS) \\ BD ≈ \leerfeld[m]
\steil Im Dreieck $ABC$ ist $AD$ die Höhe auf $BC$ ($D$ zwischen $B$ und $C$), $AD = 380$ m. $\angle ACD = 63^\circ$, $\angle CAB = 71^\circ$. Berechne $BD$. \hfill (P10 2016 OS) \\ BD ≈ \leerfeld[m]
\teil Im Drachen $ABCD$ schneiden sich die Diagonalen in $E$ rechtwinklig. Bekannt sind $AD$ und der Winkel $\delta$ bei $D$ im Dreieck $AED$. Welche Größen brauchst du, und in welcher Reihenfolge rechnest du, um den Flächeninhalt von $AED$ zu bestimmen? \hfill (P10 2015 OS)
\teil Im Parallelogramm $ABCD$ ist $DF$ die Höhe auf $AB$ ($F$ auf $AB$). Bekannt sind $AD$ und der Winkel $\alpha$ bei $A$. Welche Größen brauchst du, und in welcher Reihenfolge rechnest du, um den Flächeninhalt des Dreiecks $AFD$ zu bestimmen? \hfill (P10 2015 OS)
\end{teile}
\end{aufgabe}

% Anlauf (ohne Prüfkennung)
\begin{aufgabe}{Sinussatz – Anlauf}
\begin{teile}
\teil Im Dreieck $ABC$ ist bei $B$ ein rechter Winkel markiert, $a$ und $\alpha$ sind gegeben. Womit rechnest du $b$? \\ \kreuz{sin, cos, tan} \kreuz{Sinussatz}
\teil Dreieck $ABC$: $a = 7$ cm, $\alpha = 68^\circ$, $\beta = 51^\circ$. Wie lang ist $b$? b ≈ \leerfeld[cm]
\teil Dreieck $PQR$: $PQ = 8$ cm, Winkel bei $R$: $69^\circ$, Winkel bei $Q$: $53^\circ$. Wie lang ist $PR$? PR ≈ \leerfeld[cm]
\end{teile}
\end{aufgabe}

\begin{aufgabe}{Sinussatz – Prüfungsaufgaben}
\begin{teile}
\steil Im Drachen $ABCD$ gilt im Dreieck $ABD$: $AD = 4{,}6$ cm, der Winkel bei $A$ ist $118^\circ$, der Winkel bei $B$ ist $33^\circ$. Berechne die Diagonale $BD$. \hfill (P10 2015 OS) \\ BD ≈ \leerfeld[cm]
\steil Wanderung: Hütte $H$, Gipfel $G$, See $S$. $HG = 3\,200$ m, der Winkel bei $H$ ist $57^\circ$, der Winkel bei $S$ ist $72^\circ$. Wie lang ist der Weg vom Gipfel zum See und weiter zum Parkplatz, der $1\,800$ m hinter dem See liegt? \hfill (P10 2014 OS) \\ \leerfeld[km]
\end{teile}
\end{aufgabe}

\begin{aufgabe}{Sinussatz – Prüfungsaufgaben – weiter}
\begin{teile}
\steil Eine Seilbahn führt von $A$ über die Stütze $B$ nach $C$. $AB = 520$ m, der Winkel bei $A$ ist $31^\circ$, der Winkel $ABC$ ist $113^\circ$. Berechne die Länge der Strecke $BC$. \hfill (P10 2024 OS) \\ BC ≈ \leerfeld[m]
\steil Eine Seilbahn führt von $A$ über die Stütze $B$ nach $C$. $AB = 290$ m, der Winkel bei $A$ ist $27^\circ$, der Winkel $ABC$ ist $121^\circ$. Berechne die Länge der Strecke $BC$. \hfill (P10 2024 OS) \\ BC ≈ \leerfeld[m]
\steil Eine Rampe $y$ führt auf eine Treppe. Im Dreieck aus Rampe, waagerechter Strecke ($220$ cm) und der Verbindung Treppenfuß–obere Stufenkante liegt am Rampenfuß der Winkel $6^\circ$ und am Treppenfuß der Winkel $137^\circ$ (gegenüber $y$). Berechne $y$. \hfill (P10 2025 OS) \\ y ≈ \leerfeld[cm]
\steil Eine Rampe $y$ führt auf eine Treppe. Im Dreieck aus Rampe, waagerechter Strecke ($190$ cm) und der Verbindung Treppenfuß–obere Stufenkante liegt am Rampenfuß der Winkel $4^\circ$ und am Treppenfuß der Winkel $148^\circ$ (gegenüber $y$). Berechne $y$. \hfill (P10 2025 OS) \\ y ≈ \leerfeld[cm]
\steil Von einer Plattform aus werden zwei Sichtlinien gepeilt. Im Dreieck $ACD$ ist $CD = 9$ m, der Winkel bei $A$ zwischen den Sichtlinien beträgt $4^\circ$, der Winkel bei $D$ $62^\circ$. Zeige, dass $AC \approx 114$ m ist. \hfill (P10 2018 OS)
\steil Von einer Plattform aus werden zwei Sichtlinien gepeilt. Im Dreieck $ACD$ ist $CD = 8$ m, der Winkel bei $A$ zwischen den Sichtlinien beträgt $6^\circ$, der Winkel bei $D$ $51^\circ$. Zeige, dass $AC \approx 59{,}5$ m ist. \hfill (P10 2018 OS)
\steil Ein Grundstück $ABCD$ hat bei $B$ einen rechten Winkel. $AC = 7$ m, $\angle BAC = 59^\circ$, $\angle BCD = 88^\circ$, $\angle ADC = 78^\circ$. Berechne den Abstand $AD$. \hfill (P10 2019 OS) \\ AD ≈ \leerfeld[m]
\steil Ein Grundstück $ABCD$ hat bei $B$ einen rechten Winkel. $AC = 10$ m, $\angle BAC = 53^\circ$, $\angle BCD = 96^\circ$, $\angle ADC = 69^\circ$. Berechne den Abstand $AD$. \hfill (P10 2019 OS) \\ AD ≈ \leerfeld[m]
\steil Dreieck $DEF$: $FE = 9$ m, Winkel bei $F$ $118^\circ$, Winkel bei $E$ $23^\circ$. $P$ liegt auf $DE$ mit $PE = 8$ m. Wie lang ist $DP$? \hfill (P10 2020 OS) \\ DP ≈ \leerfeld[m]
\steil Dreieck $DEF$: $FE = 16$ m, Winkel bei $F$ $107^\circ$, Winkel bei $E$ $27^\circ$. $P$ liegt auf $DE$ mit $PE = 17$ m. Wie lang ist $DP$? \hfill (P10 2020 OS) \\ DP ≈ \leerfeld[m]
\steil Viereck $ABCD$ mit rechtem Winkel bei $A$ und der Diagonale $BD$. $\angle ADB = 48^\circ$, $\angle BDC = 57^\circ$, $\angle ABC = 104^\circ$, $DC = 1\,300$ m. Stimmt die Aussage „$BC$ ist genauso lang wie $DC$“? \hfill (P10 2021 OS)
\steil Viereck $ABCD$ mit rechtem Winkel bei $A$ und der Diagonale $BD$. $\angle ADB = 51^\circ$, $\angle BDC = 63^\circ$, $\angle ABC = 112^\circ$, $DC = 1\,760$ m. Stimmt die Aussage „$BC$ ist genauso lang wie $DC$“? \hfill (P10 2021 OS)
\end{teile}
\end{aufgabe}
